\section{Stopping and the remaining dimension budget}\label{s:stopping}

Each block estimate below loses two dimensions. Before iterating it,
we must know that enough dimensions remain. Terminal state counting
gives this fact on any putative short run. A separate full-data
argument covers constant accuracy. Both reductions count individual
samples and allow early stopping.

Throughout this section \(W=2^M\), and success is averaged over an
independent uniform signal. Logarithms denoted by \(\log_2\) are to
base two.

\begin{lemma}[Terminal count and padding]\label{s:terminal-count}
For \(d\ge3\), a learner of horizon \(T\) satisfies
\begin{align}
 \Prob(\mathrm{success})
 &\le (T+1)W(c_d/2)\epsilon^{d-1},\label{s:output-count-sharp}\\
 \Prob(\mathrm{success})
 &\le (T+1)W(2\epsilon)^{d-1}.\label{s:output-count-angular}
\end{align}
The learner can be padded to any prescribed later sample layer, with
the same success law and width at most
\begin{equation}\label{s:padded-width}
 W_*=(T+2)W.
\end{equation}
\end{lemma}
\begin{proof}
Fix a seed realization covered by Lemma~\ref{s:borel} and use its
Borel version. Its initial-state law is independent of the signal.
For each terminal index-state pair \((t,v)\), let \(q_{t,v}(s)\)
be its probability of being reached at signal \(s\), and let
\(\pi_{t,v}\) be its target-independent output law. There are at most
\((T+1)W\) such pairs. Replacing \(q_{t,v}(s)\) by one, then integrating
\(\pi_{t,v}\) over the uniform signal, bounds that pair's contribution
by a single angular cap mass. Angular success implies chordal error
at most \(\epsilon\), so \eqref{s:sharp-cap} gives
\eqref{s:output-count-sharp}; \eqref{s:angular-cap} gives
\eqref{s:output-count-angular}. The same numerical
bound holds for almost every fixed seed, so Lemma~\ref{s:borel}
permits averaging the original conditional success probabilities.

To pad, keep every halted state together with its original stopping
index as an absorbing state, and ignore later samples. At the final
layer use its original output law. At any layer there are at most
\(W\) active states and \((T+1)W\) halted index-state labels,
giving \eqref{s:padded-width}. This includes a stop at index zero.
\end{proof}

The terminal count also identifies the regime where the output
restriction alone is decisive. At success probability at least
\(2/3\), \eqref{s:output-count-sharp} gives
\[
 T+1\ge \frac{4}{3c_d}\,2^{-M}\epsilon^{-(d-1)}.
\]
For a fixed \(\alpha>0\), accuracy \(\epsilon=2^{-\alpha d}\)
therefore requires \(T\ge2^{(\alpha-o(1))d^2}\) when
\(M=o(d^2)\). The block arguments provide an additional obstruction beyond terminal
output counting. To apply them, we next extract the
remaining dimension budget from the same terminal count.

\begin{lemma}[Short-run dimension budget]\label{s:short-run}
Fix constants \(c>0\) and \(\alpha>0\), and put
\(L=\log_2(1/\epsilon)\). Suppose \(M(d)=o(d^2)\). On any sequence
of dimensions where uniform-prior success is at least \(\alpha\)
and \(T\le cdL\), one has
\begin{equation}\label{s:short-run-conclusions}
 L=o(d),\qquad T=o(d^2),\qquad \log_2 W_*=o(d^2).
\end{equation}
Consequently, for any block length \(k\ge c_0d\) with fixed \(c_0>0\),
\(\lceil T/k\rceil=o(d)\) on that sequence.
\end{lemma}
\begin{proof}
Since \(L\ge1\), \(T+1\le(c+1)dL\) and
\(\log_2 L\le L\). Taking logarithms in
\eqref{s:output-count-sharp} and using \(c_d\le\sqrt d\) gives the
estimate
\begin{equation}\label{s:sharp-short-run}
 (d-2)L
 \le M+\frac32\log_2d+\log_2\frac{c+1}{2\alpha}.
\end{equation}
This proves \(L=o(d)\), uniformly in the accuracy and the rules on
the indicated dimensions. The assumed sample bound gives
\(T=o(d^2)\); then \(\log_2 W_*=M+\log_2(T+2)=o(d^2)\).
Dividing \(T\) by \(k\ge c_0d\) proves the block-count assertion.

The short-run hypothesis is essential: the conclusion
\(L=o(d)\) applies only on the indicated dimensions.
\end{proof}

\begin{lemma}[Full-data residual bound]\label{s:full-data}
Let \(0\le T\le d-3\). Give an estimator all \(T\) Gaussian rows and
labels for a uniform signal, and allow any unit output depending on
these data and independent randomness. Put \(D=d-T\). For every
\(0<r_0<1\), its angular success probability is at most
\begin{equation}\label{s:full-data-bound}
 \frac{T}{d(1-r_0^2)}
 +\min\left\{1,\ \frac{c_D}{2}(2\epsilon/r_0)^{D-1},\
                    (8\epsilon/r_0)^{D-1}\right\}.
\end{equation}
\end{lemma}
\begin{proof}
For \(T\ge1\), the row span \(R_X\) has dimension \(T\) almost surely.
Given the full rows and labels, the row-space projection \(p=P_{R_X}S\)
is known, because the row matrix is injective on its row span.
The residual law is
\[
 S=p+rU,\qquad r=\sqrt{1-\norm p^2},\qquad
 U\sim\sigma_{R_X^\perp}.
\]
Conditional on the data, the residual direction is uniform.
At a fixed row span, spherical symmetry gives
\(\E\norm p^2=T/d\). Markov's inequality yields
\(\Prob(r<r_0)\le T/[d(1-r_0^2)]\).
On \(r\ge r_0\), successful residual directions lie in a chordal
cap of radius \(2\epsilon/r_0\), by the recentering argument in
Lemma~\ref{s:caps}. The sharp and affine terminal cap bounds give the
three alternatives in the minimum. The output randomization may be
averaged after conditioning on the full data. When \(T=0\), take
\(R_X=\{0\}\), \(p=0\), and \(r=1\); the same argument uses the
original full sphere.
\end{proof}

For \(T\le d/4\), take \(r_0=1/2\) in
\eqref{s:full-data-bound}. Since \(D\ge3d/4\) and
\(2\epsilon/r_0\le0.4\),
\[
 \Prob(\mathrm{success})
 \le\frac13+(c_D/2)0.4^{D-1}=\frac13+o(1)
\]
uniformly in this range of \(T\). Thus uniform-prior success at least
\(1/2\) requires
\begin{equation}\label{s:quarter-endpoint}
 T>d/4
\end{equation}
for all sufficiently large \(d\).



These full-data estimates apply to a learner that stopped earlier,
because revealing all pre-generated rows and labels only gives it
more information.
